\section{Conclusion}\label{sec:conclusion}

This paper set out to answer a planner's question: if climate-AI capacity can be built in only some
countries, where should it go first when the welfare of the poor counts for more? We developed
equity-weighted doubly robust policy learning to answer it. The method combines exogenous
temperature shocks, orthogonal doubly robust scores for the effect of capability on shock
sensitivity, Atkinson welfare weights over income deciles, and interpretable allocation rules
learned under a budget. We established its regret properties and showed in simulations that it
outperforms the interaction regressions common in the applied literature by a wide margin. Equity
weighting matters most when inequality aversion is high, and orthogonalisation matters for
reliable targeting.

Applied to cereal yields in 145 countries, the method finds strong heat damages but no detectable
resilience dividend from climate-AI research capacity. The estimates rule out effects larger than
about two-thirds of the damage. The allocation most often used in practice, building on existing
AI strength, delivers essentially no equity-weighted benefit, while rules that prioritise the poorest and the most
agriculturally dependent countries have higher, though imprecise, value. Inequality aversion changes the recommended
allocation only when it is high, at which point Viet Nam and most of its ASEAN neighbours enter the
priority set.

The broader lesson for research on AI, climate change and inclusive development is methodological
as well as substantive. Questions about who benefits are questions about targeting under
heterogeneity, and they call for methods that estimate allocation rules and their welfare
consequences directly. They also call for better data. The binding constraints in our application
are the absence of cross-country measures of AI deployment and the limited overlap between countries
with and without capability. Programme-level and subnational data on AI-enabled adaptation services
would allow the same method to deliver much sharper answers. The code is released with the paper
for that purpose.
